\section{Introduction and conventions}\label{sec:introduction}

The Ball--Evans approximation problem concerns the compatibility of
Sobolev approximation with global injectivity. A deformation in nonlinear
elasticity is represented by a map between spatial domains; injectivity
expresses noninterpenetration, whereas its weak first derivative records
the local deformation. Smooth approximation is elementary if injectivity
is discarded, but convolution need not preserve it. The question is
whether the two requirements can be met simultaneously. Ball's
formulation explicitly attributes the question to Evans
\cite[Section~3, p.~8]{Ball2001}; his later discussion and Hencl's survey
explain the surrounding variational questions
\cite{Ball2010,Hencl2025}.

Here a \emph{domain} is a nonempty connected open set. For domains
$\Omega,\Lambda\subset\R^3$ and $1\le p<\infty$, a Sobolev
homeomorphism is a homeomorphism $f:\Omega\to\Lambda$ whose components
and weak first derivatives belong to $L^p(\Omega)$. A smooth
diffeomorphism has a smooth inverse as well as a smooth forward map.
The fixed-target version of the approximation problem asks for
diffeomorphisms $f_j:\Omega\to\Lambda$ such that both $f_j-f$ and
$Df_j-Df$ tend to zero in $L^p(\Omega)$. In particular, every
approximating map, not just the limit, must have image exactly
$\Lambda$.

We resolve this problem for the low and critical exponents in dimension
three.

\begin{theorem}\label{main:theorem}
Let $\Omega,\Lambda\subset\R^3$ be nonempty bounded connected open sets, let $1\le p\le2$, and let $f:\Omega\to\Lambda$ be a homeomorphism in $W^{1,p}(\Omega;\R^3)$. There exist $C^\infty$ diffeomorphisms $f_j:\Omega\to\Lambda$, each belonging to $W^{1,p}(\Omega;\R^3)$, such that
\begin{equation}\label{main:convergence}
\int_\Omega\bigl( |f_j-f|^p+|Df_j-Df|^p\bigr)\,dx\longrightarrow0.
\end{equation}
\end{theorem}

Both the map and its inverse in the conclusion are smooth on their open domains. No smoothness at the boundary, boundary extension of $f$, or regularity of either boundary is assumed. Neither inverse Sobolev regularity nor a nonvanishing Jacobian is an additional hypothesis. The target $\Lambda$ is fixed throughout.

\subsection{Prior work and the three-dimensional difficulty}

In the plane, Iwaniec, Kovalev, and Onninen proved strong diffeomorphic
approximation for $1<p<\infty$
\cite[Theorem~1.1]{IwaniecKovalevOnninen2011}. Hencl and Pratelli
established the endpoint $p=1$, including approximation by locally
finite piecewise affine homeomorphisms
\cite[Theorem~1.1]{HenclPratelli2018}. These theorems preserve the
original target and do not require Sobolev regularity of the inverse.
They show that the global topological constraint can be compatible
with strong derivative control throughout the planar exponent range.
The endpoint is a distinct achievement: strict convexity of the
$L^p$ norm is unavailable at $p=1$.

The planar development also distinguishes two mechanisms. The
$p=2$ precursor of Iwaniec, Kovalev, and Onninen used harmonic
replacements and smoothing \cite{IwaniecKovalevOnninen2012}; their
general $p>1$ result used coordinatewise $p$-harmonic replacements.
Hencl and Pratelli's endpoint proof instead used geometric extension
estimates and an adapted grid. Quantitative approximation of
bi-Lipschitz planar maps by Daneri and Pratelli
\cite{DaneriPratelli2014}, and approximation of planar bi-Sobolev
$W^{1,1}$ maps by Pratelli \cite{Pratelli2017}, additionally control
inverse derivatives under their respective stronger hypotheses.
Smooth invertibility of each approximant in Theorem~\ref{main:theorem}
is not a claim of Sobolev convergence of the inverses.

Higher dimensions exhibit a genuine obstruction. Hencl and Vejnar
initiated the counterexamples at $p=1$ \cite{HenclVejnar2016}.
Campbell, Hencl, and Tengvall corrected a gap in that construction and
extended the range to $1\le p<\lfloor n/2\rfloor$ in dimensions
$n\ge4$, with Jacobian determinants of both signs on sets of positive
measure \cite{CampbellHenclTengvall2018}. A sequence of diffeomorphisms cannot
reproduce this sign pattern by strong convergence of its derivatives.
In dimension three, Hencl and Mal\'y's orientation theorem instead
gives $\det Df\ge0$ almost everywhere for a sense-preserving
$W^{1,1}_{\mathrm{loc}}$ homeomorphism
\cite[Theorem~1.1]{HenclMaly2010}. This sign conclusion does not
assert a positive Jacobian and does not furnish an approximation:
rank-deficient derivatives still require separate constructions.

Smoothing a homeomorphism that is already piecewise affine is an
important intermediate problem. Mora-Corral and Pratelli obtained
fixed-image diffeomorphic smoothing in the plane, with control of the
forward map and its inverse \cite{MoraCorralPratelli2014}.
Campbell and Soudsk\'y proved $C^1$ homeomorphic approximation of
piecewise affine maps in arbitrary dimension
\cite{CampbellSoudsky2021}; their conclusion does not require the
inverse to be smooth. Campbell, D'Onofrio, and V\'itek subsequently
proved diffeomorphic smoothing of locally finite piecewise affine
homeomorphisms in dimensions three and four, including forward and
inverse derivative-error control
\cite[Theorem~A]{CampbellDOnofrioVitek2026}. These results identify
the importance of distinguishing a smooth injective map from a
diffeomorphism. They do not construct a strong piecewise affine
approximation of an arbitrary Sobolev homeomorphism.

The present proof separates that approximation problem from smoothing.
It first constructs an onto locally bi-Lipschitz map with small strong
Sobolev error, and then proves a smoothing theorem with arbitrarily
fine pointwise control. Three-dimensional PL topology, developed by
Moise and Bing and used here in Hamilton's relative formulation,
provides qualitative replacements
\cite{Moise1952Approximation,Moise1952Triangulation,Bing1959,Hamilton1976}.
It supplies no derivative estimate. The energy bounds must therefore
come from the parametrizations and measurements used below, rather
than from the mere existence of a topological replacement.

The critical exponent $p=2$ has additional analytic structure.
Cs\"ornyei, Hencl, and Mal\'y proved that a three-dimensional
$W^{1,2}_{\mathrm{loc}}$ homeomorphism has an inverse of locally
bounded variation, and that almost every parallel plane has the
two-dimensional Lusin property
\cite[Theorems~1.1 and~1.3]{CsornyeiHenclMaly2010}. We use these two
statements, respectively, to select reverse rays and to measure
generic image surfaces. Neither is an assumption on the given inverse.
The further analytic ingredients include Whitney's extension theorem
for compatible first jets \cite{Whitney1934} and conformal
parametrization of metric disks \cite{AhlforsBers1960,McMullen2023}.
Their roles and the accompanying local estimates are made explicit
at the point of use.

The complementary range $2<p<\infty$ is treated by a separate
construction in \cite[Theorem~1.1]{OpenAIHigh2026}. The present paper does
not invoke that range theorem or any of its proof modules. The smoothing
theorem proved here holds for all finite exponents; the companion also
includes a complete proof in its Appendix~A.

\subsection{Measurements, local models, and the proof strategy}

The first objective is to replace $f$ by a locally bi-Lipschitz
homeomorphism while estimating derivatives without a quantitative
topological extension theorem. We measure the images of a locally
finite collection of source curves, and also of source surfaces when
$p=2$. A target triangulation gives complementary-dimensional dual
pieces inside its tetrahedra, defined by ties among their largest
barycentric coordinates. Intersections with these pieces determine a
\emph{budget}: each curve
crossing is charged the length of a corresponding target edge, and
each surface crossing is charged the area of a corresponding target
face, with multiplicity. Generic sampling bounds these sums by the
original trace integrals. Local topological changes make the counted
crossings regular; a face-preserving target reparametrization then
concentrates the image length or area near them. Thus
Proposition~\ref{lt:budget-transfer} converts intersection counts into
actual measurements for an onto locally bi-Lipschitz replacement.
All protected local data are retained in buffered regions that are
excluded from the measurements.

Measurements alone give an upper energy bound, not proximity to
$Df$. We therefore retain finitely many disjoint compact sets on
which the values and first derivatives of $f$ are controlled and
the derivative has a fixed positive rank. The discarded derivative
integral is small; rank-zero points carry no such energy. The
construction on these compact sets depends on the rank.
\begin{enumerate}
\item At full rank, compatible first jets extend to local
$C^1$ diffeomorphisms. Radial source changes enlarge an exact germ,
meaning agreement with such a model near the block center, into most
of a small block, so that the map recovers the prescribed
jet there. The surrounding thin region is filled using the trace
budgets. At $p=2$, reverse-ray selection and angular attenuation
control the surface cost of this radial construction
(Lemma~\ref{la:radial-blocks}).
\item At rank one or two, source coordinates split into the
nonzero directions and the kernel of a nearby fixed matrix.
Compressing the kernel reduces the controlled part of a volume block
to a line or plane section. Nearly minimal curve length gives the
correct derivative on a line. At rank two and $p=2$, the sharp
boundary-parametrized disk estimate gives the analogous conclusion
from nearly minimal area (Lemma~\ref{lt:disk}). At rank two and
$p<2$, a further radial construction inside the plane section
provides the needed approximation.
\item Uncontrolled cells are reparametrized relative to their
already chosen boundary values. Below a cell's dimension, radial
collapse bounds its energy by the boundary energy
(Lemma~\ref{lt:collapse}). This fills faces when $p<2$ and volumes
throughout $1\le p\le2$; the critical faces instead use the disk
estimate. The resulting maps preserve the old cell images and
agree on shared faces (Lemma~\ref{la:cell-completion}).
\end{enumerate}

The order of choices is part of the argument. Compact jet bounds and
all finite multiplicative constants are fixed before the excess
neighborhoods and trace errors are made small against them. The
critical radial construction separates the regions whose area can
be suppressed after replacement from those whose area must be measured
in advance; Figure~\ref{la:radial-bands} displays these regions at the
point where their construction is explained. Outside the
retained compact configurations, universal mesh constants multiply
only the small original energy tail. Summable local errors then
yield a global $W^{1,p}$ approximation, not just a locally Sobolev
map. For $p=1$, the line argument uses an elementary squared
length-defect estimate followed by H\"older's inequality, rather
than strict convexity of $L^1$.

Finally, locally bi-Lipschitz maps are approximated strongly by PL
maps using finite local model libraries. Their derivative bounds are
fixed before reducing the exceptional mesh volume. PL smoothing
then uses small edge and vertex neighborhoods with summable costs.
A fine value tolerance makes the smooth map properly homotopic to
the original homeomorphism inside the target; its positive local
degrees and degree one force bijectivity. This proves the fixed
image for each individual approximant, independently of passage
to the limit.

Section~\ref{sec:smoothing} proves this smoothing theorem first.
Section~\ref{sec:low-tools} develops the disk and collapse estimates,
relative topological replacements, and sampled trace budgets.
Section~\ref{sec:low-assembly} constructs the rank-sensitive blocks
and proves their global strong approximation estimate.
Section~\ref{sec:completion} concludes the diffeomorphic
approximation theorem on the original target.

\subsection{Conventions}

Unless stated otherwise, all domains, closures, local finiteness, and compactness are taken relative to the indicated open manifold. A compact set in an open domain has positive distance from its Euclidean complement. A PL map or triangulation is locally finite. A locally bi-Lipschitz homeomorphism has finite Lipschitz constants in both directions on sufficiently small neighborhoods; no global bound is implicit.

The symbol $|A|$ denotes the Frobenius norm of a matrix, and $\|A\|_{\op}$ its operator norm. We use the Frobenius norm in sharp gradient estimates. Equivalent norms give the same strong convergence in Theorem~\ref{main:theorem}. The symbols $\cL^m$ and $\cH^m$ denote Lebesgue and Hausdorff measure. Integrals on parameterized traces include their stated multiplicities. The differential along an $m$-dimensional parameter space is denoted $D_m$ when it must be distinguished from the full differential.

A \emph{fine value tolerance} on a domain $U$ is a positive continuous function on $U$. To approximate finely means to meet a prescribed such tolerance pointwise. Whenever inverse control is needed on a compact localization, it is included in the prescribed tests. Locally finite constructions permit tolerances tending to zero toward the ends of the open domains; no positive global lower bound is assumed.

Constants called \emph{absolute} depend only on dimension and on fixed universal reference shapes. A constant $C_p$ may also depend on the fixed exponent. Other dependencies are specified when the constant is introduced. A finite constant that depends on a chosen compact collection of jets or charts is fixed before an error is required to be small against it.

A homeomorphism between connected oriented domains has a constant topological orientation. If necessary, reflect the target, perform the construction for the orientation-preserving map, and undo the reflection. This convention imposes no extra hypothesis on the Sobolev differential.

\begin{lemma}[Local tolerances]\label{pre:local-tolerances}
Let $U\subset\R^n$ be open and let $a:U\to(0,\infty)$ be continuous. For every $L>0$ there is a positive $L$-Lipschitz function $d$ on $U$ such that
\[
 d(x)\le \min\{a(x),1,L\dist(x,\R^n\setminus U)\}.
\]
There are also positive smooth minorants satisfying prescribed positive continuous upper bounds on their values and first derivatives. One may additionally impose positive constant bounds on a locally finite family of compact regional tests. Finite collections of such requirements can be combined, and error budgets on a countable compact exhaustion can be made summable.
\end{lemma}
\begin{proof}
Extend $e(x)=\min\{a(x),1,L\dist(x,\R^n\setminus U)\}$ by zero outside $U$, and put
\[
 d(x)=\inf_{y\in\R^n}\{e(y)+L|x-y|\}.
\]
This function is $L$-Lipschitz and is bounded above by $e$. It is positive at an interior point $x$: on a small closed ball about $x$, $e$ has a positive minimum, and outside that ball the distance term has a positive lower bound.

For a smooth minorant take a smooth locally finite partition of unity $\{\chi_i\}$ with compact supports in $U$ and put $b=\sum_i c_i\chi_i$, with positive constants $c_i$. Given positive continuous bounds $v,w$, choose
\[
 c_i\le 2^{-i-2}\min\left\{
 \inf_{\supp\chi_i}v,
 \frac{\inf_{\supp\chi_i}w}{1+\|D\chi_i\|_\infty}
 \right\}.
\]
Then $b>0$, $b\le v$, and $|Db|\le w$. Replace a finite list of requirements by their minimum. For countably many regional tests use a locally finite compact cover with compact enlargements, refine the bounds on each intersecting support, and allocate total errors by a convergent positive series. Only finitely many regional requirements meet a given compact support.
\end{proof}

We will also use the following elementary global observation. If $H:U\to V$ is a homeomorphism and a continuous map $G:U\to V$ satisfies
\[
 |G(x)-H(x)|<\tfrac12\dist(H(x),\R^n\setminus V),
\]
then the straight homotopy stays in $V$. When $V$ is bounded, it is a proper homotopy: the preimage of a compact target set is contained, uniformly in the homotopy parameter, in the $H$-preimage of a compact set a positive distance from $\partial V$. This observation will be used together with local injectivity and degree, not as a substitute for local injectivity.
